% Insertion-ready fragment for the existing quaternionic module.
% It uses the already defined original P,R,u,C,X,B_0,Q_latt,S,C_2,C_3,
% Z_3,Omega_3 and the monodromy M_3 from q444:period-gram.
% No canonical file was edited by the independent checking agent.

\subsection*{All positive period-metric fibres and the cost of rank loss}

The following calculation keeps the periods in
Proposition~\ref{q444:period-gram} fixed.  In particular
$C_3(t)=SC_2(t)Q_{\rm latt}$ still uses the displayed $S$, not the
different right inverse $R=P^*C^{-1}$.  Choose a punctured cusp on which
the retained correction is bounded, and fix a finite constant $K\geq0$
such that $\|C_3(t)\|\leq K$ there.  Such a constant exists because
$\|C_3(t)\|\leq\|S\|\|C_2(t)\|\|Q_{\rm latt}\|$.
All norms below are the Euclidean operator norms.  Write
\[
 t=e^{2\pi i s},\qquad s=x+iy,\qquad
 y=-\frac{\log|t|}{2\pi},
\]
and restrict to the nonempty smaller cusp $y>\max\{K,1\}$.
This changes no period or coefficient.  It identifies a domain where
the inverse used below is proved to exist.

\begin{theorem}[The complete positive metric fibre]
\label{q444:metric-fibre}
Retain
\[
 A_{\rm per}(s)=xI_3+\operatorname{Re}C_3(t),\qquad
 Y_{\rm per}(t)=yI_3+\operatorname{Im}C_3(t),
\]
and define the real period map and its inverse by
\[
 \mathcal N(s)=
 \begin{pmatrix}I_3&A_{\rm per}(s)\\0&Y_{\rm per}(t)\end{pmatrix},
 \qquad
 \mathcal N(s)^{-1}=
 \begin{pmatrix}
 I_3&-A_{\rm per}(s)Y_{\rm per}(t)^{-1}\\
 0&Y_{\rm per}(t)^{-1}
 \end{pmatrix}.                                             \tag{MF.1}
\]
For every prescribed Hermitian matrix $H=E+iF\succ0$, where
$E=E^{\mathsf T}$ and $F=-F^{\mathsf T}$ are real, all real symmetric
positive-definite lattice matrices giving the unchanged value
$\Omega_3W_3\Omega_3^*=H$ are exactly
\[
 W_3=\mathcal N^{-1}
 \left[
  \frac12\begin{pmatrix}E&-F\\F&E\end{pmatrix}
  +\begin{pmatrix}U&V\\V&-U\end{pmatrix}
 \right]\mathcal N^{-\mathsf T},                             \tag{MF.2}
\]
where $U,V$ are real symmetric $3\times3$ matrices and the bracket is
positive definite.  This is a nonempty open convex twelve-real-dimensional
fibre.  An explicit description of its entire parameter domain is
\[
 E_+=E/2+U\succ0,\qquad
 E/2-U-(-F/2+V)^{\mathsf T}E_+^{-1}(-F/2+V)\succ0.           \tag{MF.3}
\]
In particular $U=V=0$ always gives a positive lift.

Every smooth map $H_3(t)\succ0$ on the smaller punctured cusp has a
smooth positive lift (MF.2) with
\[
 W_3(s+1)=M_3^{-1}W_3(s)M_3^{-\mathsf T}.
                                                               \tag{MF.4}
\]
The full set of positive observations $G(t)=PH_3(t)P^*$ obtainable
from such metrics is exactly the set of all smooth
$\operatorname{Herm}_2(\mathbb C)_{\succ0}$-valued maps on that cusp.
For fixed $G\succ0$, all the intermediate positive $H$ are
\[
 H=X\begin{pmatrix}
       G&z\\ z^*&z^*G^{-1}z+\tau
     \end{pmatrix}X^*,\qquad
 z\in\mathbb C^2,\quad \tau>0.                                \tag{MF.5}
\]
The positive real lattice-metric fibre of $G$ has real dimension
seventeen.
\end{theorem}

\begin{proof}
For every real vector $v$,
\[
 \|Y_{\rm per}v\|
 \geq y\|v\|-\|\operatorname{Im}C_3\,v\|
 \geq(y-K)\|v\|.
\]
The inequality $\|\operatorname{Im}C_3\|\leq\|C_3\|$ follows from
$\operatorname{Im}C_3=(C_3-\bar C_3)/(2i)$ and
$\|\bar C_3\|=\|C_3\|$.  Since $y>K$, $Y_{\rm per}$ has zero kernel
and is an invertible square real matrix.  Block multiplication now proves
both inverse identities in (MF.1).  This is the original real period map:
if $a,b\in\mathbb R^3$, then
\[
 a+Z_3b=(a+A_{\rm per}b)+iY_{\rm per}b.
\]
Thus the inverse takes $v+iw$ to
$(v-A_{\rm per}Y_{\rm per}^{-1}w,\;Y_{\rm per}^{-1}w)$.
No correction in $C_3$ has been removed.

Put $\mathcal J=(I_3\;\;iI_3)$.  The exact factorization is
$\Omega_3=\mathcal J\mathcal N$.  Congruence gives a bijection between
real positive $W_3$ and real positive
$T=\mathcal NW_3\mathcal N^{\mathsf T}$.  For
\[
 T=\begin{pmatrix}E_0&F_0\\F_0^{\mathsf T}&K_0\end{pmatrix}
\]
one obtains by multiplication
\[
 \mathcal JT\mathcal J^*
   =E_0+K_0+i(F_0^{\mathsf T}-F_0).                            \tag{MF.6}
\]
Consequently this equals $E+iF$ precisely when
$E_0=E/2+U$, $K_0=E/2-U$, and $F_0=-F/2+V$ with $U,V$ symmetric.
Indeed, the first equality in (MF.6) fixes $K_0=E-E_0$,
and $F_0^{\mathsf T}-F_0=F$ is equivalent to symmetry of $F_0+F/2$.
These recover $U,V$ uniquely, proving the full fibre formula.

For real $v,w$, with $q=v+iw$,
\[
 \begin{pmatrix}v\\w\end{pmatrix}^{\mathsf T}
 \begin{pmatrix}E&-F\\F&E\end{pmatrix}
 \begin{pmatrix}v\\w\end{pmatrix}=q^*Hq.
\]
The right side is positive for every nonzero $(v,w)$, so the first
bracket term in (MF.2) is positive definite.  This proves nonemptiness.
The kernel parameters have dimension $6+6=12$.  Positive definiteness
is an open convex condition on real symmetric matrices, so its intersection
with this affine twelve-dimensional fibre is open and convex and has
dimension twelve.  Finally, for $B_+=-F/2+V$,
\[
 \begin{pmatrix}E_+&B_+\\B_+^{\mathsf T}&E/2-U\end{pmatrix}
 =
 \begin{pmatrix}I&0\\B_+^{\mathsf T}E_+^{-1}&I\end{pmatrix}
 \begin{pmatrix}E_+&0\\0&E/2-U-B_+^{\mathsf T}E_+^{-1}B_+\end{pmatrix}
 \begin{pmatrix}I&E_+^{-1}B_+\\0&I\end{pmatrix}.
\]
The two triangular factors are transposes of each other and are invertible;
the identity proves the necessity and sufficiency of
(MF.3), including every positivity restriction on the kernel parameters.

The original correction is single-valued in $t$, so
\[
 \mathcal N(s+1)=\mathcal N(s)M_3,\qquad
 \mathcal N(s+1)^{-1}=M_3^{-1}\mathcal N(s)^{-1}.               \tag{MF.7}
\]
For any smooth prescribed $H_3(t)$, its real and imaginary parts are smooth
and periodic on the cover.  The inverse of $Y_{\rm per}$ is smooth,
since its entries are its adjugate divided by its nonzero determinant.
Formula (MF.2) with $U=V=0$ is therefore smooth and satisfies (MF.4).
More generally, smooth periodic $U,V$ satisfying (MF.3) give all such
smooth lifts: conversely (MF.7) and (MF.4) make
$\mathcal NW_3\mathcal N^{\mathsf T}$ periodic, and its blocks recover
periodic $U,V$ by (MF.6).  No holomorphic condition is being imposed on
the real metric or on the prescribed Hermitian map.

If $W_3\succ0$ is real, it is also positive on nonzero complex columns:
for $p+iq$, its quadratic form is $p^{\mathsf T}W_3p+
q^{\mathsf T}W_3q$.  Since $\Omega_3^*$ is injective by the identity
block of $\Omega_3$, one has $H_3\succ0$.  Since $P^*$ is injective,
also $G=PH_3P^*\succ0$.  This proves necessity of strict positivity;
no singular $G$ occurs at an interior point for a positive metric.

For the converse, the marked inverse $X^{-1}=(P;\,u^*/3)$ makes the
upper-left block of $X^{-1}HX^{-*}$ equal to $PH P^*$.  With that block
fixed at $G\succ0$, the complete positive block matrices are
\[
 \begin{pmatrix}G&z\\z^*&\ell\end{pmatrix},\qquad
 \ell-z^*G^{-1}z=\tau>0.
\]
To prove this exact positivity statement, multiply
\[
 \begin{pmatrix}I_2&0\\z^*G^{-1}&1\end{pmatrix}
 \begin{pmatrix}G&0\\0&\tau\end{pmatrix}
 \begin{pmatrix}I_2&G^{-1}z\\0&1\end{pmatrix}.
\]
The result is the displayed block with $\ell=z^*G^{-1}z+\tau$.
The invertible triangular factors prove both directions, and congruence
by $X$ gives (MF.5).  Taking $z=0$ and any fixed $\tau=\lambda>0$
already gives the smooth lift $H_3(t)=RG(t)R^*+\lambda uu^*$.
Applying (MF.2) proves every asserted positive observation is realized.
The intermediate $H$ fibre has four real coordinates in $z$ and one in
$\tau$.  Equivalently, the real-linear metric-to-observation map has rank
four: the constructions above give a right inverse on all Hermitian
matrices before imposing positivity.  Its real symmetric domain has
dimension $6\cdot7/2=21$, so its affine kernel has dimension $21-4=17$.
The positive part is nonempty and relatively open, proving the stated
seventeen-dimensional fibre without a dimension-only surjectivity argument.
\end{proof}

The freedom above does not replace the retained periods by an identity
period matrix.  It computes the metric which compensates their full real
period map.  The correction itself has not lost information in passing
from rank two to rank three: the exact inverse on its prescribed image is
\[
 C_2=P C_3 S B_0,
\]
because $PS=I_2$ and $Q_{\rm latt}S=B_0^{-1}$.
To identify exactly what an unrestricted metric observation cannot recover,
take two retained corrections with real period maps $\mathcal N$ and
$\widetilde{\mathcal N}$ on a common sufficiently small cusp.  The
invertible real transformation
\[
 T_{\rm per}=\widetilde{\mathcal N}^{-1}\mathcal N,\qquad
 \widetilde W_3=T_{\rm per}W_3T_{\rm per}^{\mathsf T}
                                                               \tag{MF.8}
\]
preserves positivity and gives
$\widetilde\Omega_3\widetilde W_3\widetilde\Omega_3^*
=\Omega_3W_3\Omega_3^*$ by the two exact factorizations
$\Omega_3=\mathcal J\mathcal N$.
Furthermore $T_{\rm per}(s+1)=M_3^{-1}T_{\rm per}(s)M_3$, so (MF.8)
preserves the prescribed metric monodromy.  This is an explicit
correspondence between period--metric pairs with the same $H_3$ and $G$.
Thus these observations alone cannot recover the period correction when
the metric is allowed this freedom.  The constants remain in the metric
formula; they have not been set to zero.

\begin{proposition}[Relative-metric bounds obstruct projective rank loss]
\label{q444:metric-condition}
For an equivariant positive metric put
\[
 W_*(t)=M(x)W_3(s)M(x)^{\mathsf T},\qquad
 \widetilde\Omega_3(t)=(I_3\;\;iyI_3+C_3(t)).
\]
These are single-valued on the punctured cusp, and
$H_3=\widetilde\Omega_3W_*\widetilde\Omega_3^*$.
If fixed constants $0<a\leq b<\infty$ obey
$aI_6\preceq W_*(t)\preceq bI_6$, then
\[
 a\bigl(1+(y-K)^2\bigr)C
 \preceq G(t)\preceq
 b\bigl(1+(y+K)^2\bigr)C.                                    \tag{MC.1}
\]
In particular the unchanged eigenvalues satisfy
\[
 \lambda_{\min}(G)\geq a\bigl(1+(y-K)^2\bigr)\lambda_{\min}(C),
 \qquad
 \lambda_{\max}(G)\leq b\bigl(1+(y+K)^2\bigr)\lambda_{\max}(C).
                                                               \tag{MC.2}
\]
With the original $C$, $\lambda_{\min}(C)=1$ and
$\lambda_{\max}(C)=3$.  More generally, writing
$\kappa_*(t)=\lambda_{\max}(W_*)/\lambda_{\min}(W_*)$, one always has
\[
 \frac{\lambda_{\min}(G)}{\lambda_{\max}(G)}
 \geq\frac{1+(y-K)^2}{3\kappa_*(t)(1+(y+K)^2)}.               \tag{MC.3}
\]
Thus if the observed eigenvalue ratio tends to zero as $t\to0$, the
relative metric condition number tends to infinity.  Bounded relative
condition number excludes this projective rank loss even when the common
scale of the metric is not bounded above or below.
\end{proposition}

\begin{proof}
The identity $M(x+1)=M(x)M_3$, together with (MF.4), proves that $W_*$
is periodic.  Multiplication gives
$\Omega_3M(-x)=\widetilde\Omega_3$, which proves the stated Gram formula.
For every complex vector $v$,
\[
 (y-K)\|v\|\leq\|(iyI_3+C_3)^*v\|\leq(y+K)\|v\|.
\]
Both inequalities are the triangle inequality, its reverse, and
$\|C_3^*\|\leq K$.  Since
$\widetilde\Omega_3\widetilde\Omega_3^*
=I_3+(iyI_3+C_3)(iyI_3+C_3)^*$, they give
\[
 \bigl(1+(y-K)^2\bigr)I_3
 \preceq\widetilde\Omega_3\widetilde\Omega_3^*
 \preceq\bigl(1+(y+K)^2\bigr)I_3.
\]
Congruencing $aI_6\preceq W_*\preceq bI_6$ by
$\widetilde\Omega_3$, and then by $P$, proves (MC.1).
The minimum and maximum Rayleigh quotients prove (MC.2).
The characteristic polynomial of the original $C$ is
$(z-1)(z-3)$, retaining the factors $1$ and $3$.

At each point take
$a(t)=\lambda_{\min}(W_*)$ and $b(t)=\lambda_{\max}(W_*)$ in the
same inequalities.  Dividing the resulting positive scalar bounds gives
(MC.3); this division is a comparison of eigenvalues, not a rescaling or
replacement of $G$ or $W_*$.  Its period-dependent factor tends to $1$.
For fixed $a,b$, the lower bound tends to $a/(3b)>0$; for example it is
at least $a/(27b)$ once $y\geq2K$.
If the ratio on the left tends to zero, (MC.3) implies
\[
 \kappa_*(t)\geq
 \frac{1+(y-K)^2}
 {3(1+(y+K)^2)}
 \frac{\lambda_{\max}(G)}{\lambda_{\min}(G)}
 \longrightarrow+\infty .
\]
This proves the full asserted obstruction, including metrics whose
overall scale varies.
\end{proof}

\begin{proposition}[A retained-period rank-loss family and its exact cost]
\label{q444:rank-loss-metric}
Fix, without rescaling any matrix, four real constants
\[
 \alpha_{\rm obs}>0,\qquad \beta_{\rm obs}>0,\qquad
 \lambda>0,\qquad p>0,
\]
and set
\[
 \delta(t)=\beta_{\rm obs}y^{-p},\qquad
 G(t)=\begin{pmatrix}\alpha_{\rm obs}&0\\0&\delta(t)\end{pmatrix},
 \qquad H_3(t)=RG(t)R^*+\lambda uu^*.                         \tag{RL.1}
\]
Take the lift (MF.2) with $U=V=0$.  It is a smooth equivariant positive
real lattice metric at every point of the smaller punctured cusp, with
the original periods and correction unchanged.  Its observation has
\[
 \det G(t)=\alpha_{\rm obs}\beta_{\rm obs}y^{-p}
 =\alpha_{\rm obs}\beta_{\rm obs}
       \left(-\frac{\log|t|}{2\pi}\right)^{-p}\longrightarrow0,
 \qquad
 G(t)\longrightarrow\operatorname{diag}(\alpha_{\rm obs},0).
                                                               \tag{RL.2}
\]
Put $D_2(t)=\operatorname{Im}C_2(t)$ and
$Y=Y_{\rm per}(t)=yI_3+\operatorname{Im}C_3(t)$.  The relative metric
has the exact determinant
\[
 \det W_*(t)=
 \frac{\alpha_{\rm obs}^2\lambda^2\beta_{\rm obs}^2\,y^{-2p}}
      {64\,\det(Y)^2},\qquad
 \det(Y)=y\left[
 y^2+y\,\operatorname{tr}(D_2B_0^{-1})+\det D_2\right].        \tag{RL.3}
\]
In particular
\[
 \det W_*(t)\sim
 \frac{\alpha_{\rm obs}^2\lambda^2\beta_{\rm obs}^2}{64}
 y^{-2p-6}.                                                   \tag{RL.4}
\]
If $L_{\max}=\max\{2\alpha_{\rm obs}/3,\,3\lambda\}$, then
\[
 \lambda_{\min}(W_*)\sim\frac{\beta_{\rm obs}}4\,y^{-p-2},
 \qquad
 \lambda_{\max}(W_*)\longrightarrow\frac{L_{\max}}2,
 \qquad
 \kappa_*(t)\sim\frac{2L_{\max}}{\beta_{\rm obs}}\,y^{p+2}.
                                                               \tag{RL.5}
\]
Every constant and the complete correction in the exact formulas
(RL.1)--(RL.3) is retained.
\end{proposition}

\begin{proof}
Formula (MF.5) with $z=0$ and $\tau=\lambda$ gives $H_3\succ0$ and
$PH_3P^*=G$.  Formula (MF.2) proves the metric, smoothness, and monodromy
assertions.  Equation (RL.2) is the direct determinant of the displayed
matrix and uses the original logarithmic-cover factor $2\pi$.
The ratio of its least and greatest eigenvalues is eventually
$\delta/\alpha_{\rm obs}$, so this is projective rank loss in the
precise sense used in Proposition~\ref{q444:metric-condition}.

Write $A_c=\operatorname{Re}C_3$ and
\[
 \widetilde{\mathcal N}=
 \begin{pmatrix}I_3&A_c\\0&Y\end{pmatrix},\qquad
 \mathcal N=\widetilde{\mathcal N}M(x).
\]
Here $H_3$ is real symmetric.  The relative metric is exactly
\[
 W_*=\frac12\widetilde{\mathcal N}^{-1}
       \begin{pmatrix}H_3&0\\0&H_3\end{pmatrix}
       \widetilde{\mathcal N}^{-\mathsf T}
 =\frac12\begin{pmatrix}
 H_3+A_cY^{-1}H_3Y^{-\mathsf T}A_c^{\mathsf T}
       &-A_cY^{-1}H_3Y^{-\mathsf T}\\
 -Y^{-1}H_3Y^{-\mathsf T}A_c^{\mathsf T}
       &Y^{-1}H_3Y^{-\mathsf T}
 \end{pmatrix}.                                               \tag{RL.6}
\]
Since $H_3=X\operatorname{diag}(G,\lambda)X^*$ and $\det X=1$,
$\det H_3=\lambda\alpha_{\rm obs}\delta$.
Taking the real $6\times6$ determinant in (RL.6) gives
$\det W_*=(\det H_3)^2/(2^6\det Y^2)$, proving the first identity
in (RL.3).

The second identity retains the correction rather than replacing $Y$
by $yI_3$.  Indeed
$Y=yI_3+SD_2Q_{\rm latt}$ and $Q_{\rm latt}S=B_0^{-1}$.
For rectangular $U,V$, the identity
$\det(I+UV)=\det(I+VU)$ follows by eliminating either off-diagonal
block of $\left(\begin{smallmatrix}I&U\\-V&I\end{smallmatrix}\right)$
with triangular determinant-one row operations.  Applying it here gives
\[
 \det Y
 =y^3\det\left(I_2+y^{-1}D_2B_0^{-1}\right)
 =y\left[y^2+y\operatorname{tr}(D_2B_0^{-1})+\det D_2\right],
\]
because $\det B_0=1$.  If
$D_2=\left(\begin{smallmatrix}d_{11}&d_{12}\\d_{21}&d_{22}
\end{smallmatrix}\right)$, this is the complete polynomial
\[
 \det Y=y\left[y^2+y(d_{12}-d_{21})
                     +d_{11}d_{22}-d_{12}d_{21}\right].
\]
The entries of $D_2$ remain the prescribed bounded functions of $t$.
Thus $\det Y/y^3\to1$, proving (RL.4) with the displayed coefficient.

For the eigenvalue assertions, $H_3u=3\lambda u$ because $R^*u=0$.
On the two-dimensional invariant plane $u^\perp$, its trace is
\[
 \operatorname{tr}(C^{-1}G)
      =\frac23(\alpha_{\rm obs}+\delta)
\]
and the product of the two eigenvalues is
$\det H_3/(3\lambda)=\alpha_{\rm obs}\delta/3$.
The invariance of the plane follows also from Hermiticity and the
displayed $u$ eigenvector.  Its two eigenvalues are therefore exactly
\[
 h_\pm=
 \frac{\alpha_{\rm obs}+\delta
       \pm\sqrt{\alpha_{\rm obs}^2-\alpha_{\rm obs}\delta+\delta^2}}3.
                                                               \tag{RL.7}
\]
They are positive: their sum and product are positive.  Rationalizing
the numerator of $h_-$ gives
\[
 \frac{h_-}{\delta}
 =\frac{\alpha_{\rm obs}}
 {\alpha_{\rm obs}+\delta+
       \sqrt{\alpha_{\rm obs}^2-\alpha_{\rm obs}\delta+\delta^2}}
 \longrightarrow\frac12.
\]
Also $h_+\to2\alpha_{\rm obs}/3$.
Consequently $\lambda_{\min}(H_3)\sim\delta/2$ and
$\lambda_{\max}(H_3)\to L_{\max}$, with $H_3$ bounded.

The explicit blocks (RL.6), the bound
$\|Y^{-1}\|\leq(y-K)^{-1}$, and bounded $A_c,H_3$ imply
\[
 \left\|W_*-\tfrac12\operatorname{diag}(H_3,0)\right\|
 =O(y^{-2}).
\]
For real symmetric matrices the maximum eigenvalue changes by at most
the operator norm of a perturbation: apply the Rayleigh-quotient maximum
to each matrix and subtract.  Hence
$\lambda_{\max}(W_*)\to L_{\max}/2$.

To control the least eigenvalue without assuming a bounded condition
number of $H_3$, use the exact inverse
\[
 W_*^{-1}
 =2\widetilde{\mathcal N}^{\mathsf T}
      \operatorname{diag}(H_3^{-1},H_3^{-1})
      \widetilde{\mathcal N}.
\]
For $y>1$, the decomposition
$\widetilde{\mathcal N}=\operatorname{diag}(I_3,yI_3)
+\left(\begin{smallmatrix}0&A_c\\0&\operatorname{Im}C_3\end{smallmatrix}\right)$
gives $\|\widetilde{\mathcal N}\|\leq y+\sqrt2K$.
Thus
\[
 \lambda_{\max}(W_*^{-1})
 \leq\frac{2(y+\sqrt2K)^2}{\lambda_{\min}(H_3)}.
\]
Its lower-right principal block includes the positive term
$2Y^{\mathsf T}H_3^{-1}Y$.  If $v$ is a unit eigenvector of
$H_3^{-1}$ for its greatest eigenvalue, testing that block on
$Y^{-1}v/\|Y^{-1}v\|$ gives
\[
 \lambda_{\max}(W_*^{-1})
 \geq\frac{2(y-K)^2}{\lambda_{\min}(H_3)}.
\]
All denominators are positive.  Taking reciprocals yields
\[
 \frac{\lambda_{\min}(H_3)}{2(y+\sqrt2K)^2}
 \leq\lambda_{\min}(W_*)
 \leq\frac{\lambda_{\min}(H_3)}{2(y-K)^2}.                     \tag{RL.8}
\]
Both bounds are asymptotic to $\delta/(4y^2)$, proving the first
assertion of (RL.5).  Combining it with the maximum-eigenvalue limit
proves the condition-number asymptotic with its full coefficient.
\end{proof}

This is metric-choice freedom, not a selection of a physical momentum
or mass.  For the same unchanged periods the allowed choice $W_*=I_6$
instead gives exactly
\[
 H_3=I_3+(iyI_3+C_3)(iyI_3+C_3)^*,\qquad
 G=y^2C+O(y),\qquad \det G=3y^4+O(y^3),
\]
whereas (RL.1)--(RL.2) give a determinant tending to zero and retain the
same $C_2,C_3$.  These two explicit admissible choices, the complete
twelve-dimensional metric fibres, and the seventeen-dimensional observed
fibres prove that the specified period data and positivity alone do not
select an observed trajectory.  The eigenvalue ratios above only diagnose
rank loss; no field, period, metric, or parameter was replaced by a
normalized one.  Units, a metric-selection law, dynamics, and an
interpretation of the marked observation would be additional specified
data.  No extension across other finite-order overlaps or special
fillings, and no physical or quantum gap conclusion, is asserted here.
